\documentclass[11pt]{beamer}
\usetheme{CambridgeUS}
\usecolortheme{orchid}

% ============================================================
% PACOTES
% ============================================================
\usepackage[utf8]{inputenc}
\usepackage[brazil]{babel}
\usepackage[T1]{fontenc}

\usepackage{amsmath}
\usepackage{amsfonts}
\usepackage{amssymb}
\usepackage{graphicx}
\usepackage{setspace}
\usepackage{ragged2e}
\usepackage{gensymb}
\usepackage{fontawesome5}
\usepackage{hyperref}

\setbeamertemplate{caption}[numbered]
\setbeamercovered{transparent} 
\setbeamertemplate{navigation symbols}{}

% ============================================================
% DADOS — CAPA
% ============================================================
\author[Joao Victor]{%
	Joao Victor\\[0.3em]
	\small\texttt{2024007388@ifam.edu.br}}
\title[IA e Educação Inclusiva]{O uso da Inteligência Artificial na Acessibilidade
	Educacional}
%\subtitle{Do propósito à prática com ferramentas reais}
\institute[]{%
	\textsuperscript{} Graduando em Licenciatura em Matemática\\
%	\small Instituto Federal de Educação, Ciência e Tecnologia
}
\date{\today}

% ============================================================
% COMANDOS AUXILIARES
% ============================================================
\newcommand{\citabnt}[1]{%
	\vfill
	\begin{flushright}
		\tiny\textit{#1}
	\end{flushright}
}

% ============================================================
% DOCUMENTO
% ============================================================
\begin{document}
	
	% ---------- CAPA ----------
	\begin{frame}
		\titlepage
	\end{frame}
	
	% ---------- AGENDA ----------
	\begin{frame}{Agenda}
		\begin{center}
			\Large
			\textbf{POR QUÊ} $\rightarrow$ \textbf{COMO} $\rightarrow$ \textbf{O QUÊ}
		\end{center}
		\vspace{1.5em}
		\begin{columns}[T]
			\begin{column}{0.33\textwidth}
				\centering
				\textbf{\large POR QUÊ}\\[0.3em]
				\tiny Propósito
			\end{column}
			\begin{column}{0.33\textwidth}
				\centering
				\textbf{\large COMO}\\[0.3em]
			
				\tiny Prática
			\end{column}
			\begin{column}{0.33\textwidth}
				\centering
				\textbf{\large O QUÊ}\\[0.3em]
			
				\tiny Síntese
			\end{column}
		\end{columns}
		\citabnt{Adaptado de SINEK; AUGUSTO et al. (2026)}
	\end{frame}
	
	% ============================================================
	% BLOCO 1 — POR QUÊ
	% ============================================================
	\section{POR QUÊ}
	
	\begin{frame}{}
		\begin{center}
			\Huge POR QUÊ
			\vspace{1em}
			\normalsize
			O propósito que move tudo
		\end{center}
	\end{frame}
	
	\begin{frame}{A pergunta}
		\begin{center}
			\Huge
			Por que nem todos aprendem?
		\end{center}
	\end{frame}
	
	\begin{frame}{O direito}
		\begin{center}
			\Large
			\textbf{Educação é direito de todos.}\\
			\vspace{0.5em}
			Não privilégio de alguns.
		\end{center}
		\citabnt{AUGUSTO et al. (2026)}
	\end{frame}
	
	\begin{frame}{As barreiras}
		\begin{center}
			\Large \textbf{O que impede?}
		\end{center}
		\vspace{1.5em}
		\begin{columns}[T]
			\begin{column}{0.33\textwidth}
				\centering
				\faIcon{blind}\\[0.3em]
				Físicas\\ Comunicacionais
			\end{column}
			\begin{column}{0.33\textwidth}
				\centering
				\faIcon{laptop}\\[0.3em]
				Tecnológicas\\ Metodológicas
			\end{column}
			\begin{column}{0.33\textwidth}
				\centering
				\faIcon{users}\\[0.3em]
				Atitudinais
			\end{column}
		\end{columns}
		\citabnt{AUGUSTO et al. (2026)}
	\end{frame}
	
	\begin{frame}{A virada}
		\begin{center}
			\Large
			\textbf{E se a tecnologia pudesse}\\
			\textbf{remover barreiras?}
			\vspace{1em}
			\normalsize
			Não substituir. \textbf{Ampliar.}
		\end{center}
	\end{frame}
	
	\begin{frame}{O propósito}
		\begin{center}
			\Large \textbf{Incluir é garantir:}
		\end{center}
		\vspace{1.5em}
		\begin{center}
			\Large
			Acesso $\cdot$ Permanência $\cdot$ Participação $\cdot$ Aprendizagem
		\end{center}
		\citabnt{AUGUSTO et al. (2026); UNESCO (1994)}
	\end{frame}
	
	\begin{frame}{Atividade — 15 min}
		\begin{center}
			\Large \textbf{Roda de conversa}
		\end{center}
		\vspace{1.5em}
		\begin{itemize}
			\item Quem fica de fora?
			\item O que já tentamos?
			\item O que ainda falta?
		\end{itemize}
	\end{frame}
	
	% ============================================================
	% BLOCO 2 — COMO
	% ============================================================
	\section{COMO}
	
	\begin{frame}{}
		\begin{center}
			\Huge COMO
			\vspace{1em}
			\normalsize
			Ferramentas, mãos na massa e prática
		\end{center}
	\end{frame}
	
	\begin{frame}{Mapa das ferramentas}
		\begin{center}
			\Large \textbf{As IAs que vamos usar}
		\end{center}
		\vspace{1em}
		\begin{columns}[T]
			\begin{column}{0.5\textwidth}
				\begin{itemize}
					\item \textbf{ChatGPT} — assistente geral
					\item \textbf{DeepSeek} — raciocínio e código
					\item \textbf{Perplexity} — pesquisa com fontes
					\item \textbf{NotebookLM} — análise de documentos
				\end{itemize}
			\end{column}
			\begin{column}{0.5\textwidth}
				\begin{itemize}
					\item \textbf{Japiim (Bemol)} — IA do Norte
					\item \textbf{VLibras} — tradução Libras
					\item \textbf{Seeing AI} — visão computacional
					\item \textbf{Google Meet} — legendagem
				\end{itemize}
			\end{column}
		\end{columns}
		\citabnt{AUGUSTO et al. (2026)}
	\end{frame}
	
	% ----- 1. ChatGPT -----
	\begin{frame}{\faIcon{robot}~ChatGPT}
		\textbf{O que é:} assistente conversacional da OpenAI
		\vspace{0.5em}
		\begin{block}{Como usar na educação inclusiva}
			\small
			• Criar tutor personalizado por disciplina\\
			• Adaptar textos para diferentes níveis de leitura\\
			• Gerar atividades diferenciadas\\
			• Explicar conceitos em linguagem simples
		\end{block}
		\textbf{Exemplo de prompt:}
		\textit{``Adapte este texto para um aluno com dislexia, use frases curtas e linguagem simples.''}
		\citabnt{OpenAI (2024); AUGUSTO et al. (2026)}
	\end{frame}
	
	% ----- 2. DeepSeek -----
	\begin{frame}{\faIcon{brain}~DeepSeek}
		\textbf{O que é:} IA com foco em raciocínio e código
		\vspace{0.5em}
		\begin{block}{Como usar na educação inclusiva}
			\small
			• Resolver problemas matemáticos passo a passo\\
			• Criar materiais em LaTeX acessíveis\\
			• Gerar descrições alternativas para imagens\\
			• Apoiar programação para alunos com deficiência
		\end{block}
		\textbf{Diferencial:} gratuito, open-source, forte em raciocínio lógico
		\citabnt{DeepSeek (2024); AUGUSTO et al. (2026)}
	\end{frame}
	
	% ----- 3. Perplexity -----
	\begin{frame}{\faIcon{search}~Perplexity}
		\textbf{O que é:} buscador com IA que cita fontes
		\vspace{0.5em}
		\begin{block}{Como usar na educação inclusiva}
			\small
			• Pesquisar legislação sobre inclusão com fontes\\
			• Encontrar artigos científicos atualizados\\
			• Verificar informações sobre tecnologias assistivas\\
			• Montar bibliografias confiáveis
		\end{block}
		\textbf{Diferencial:} respostas com referências verificáveis
		\citabnt{Perplexity (2024); AUGUSTO et al. (2026)}
	\end{frame}
	
	% ----- 4. NotebookLM -----
	\begin{frame}{\faIcon{book}~NotebookLM}
		\textbf{O que é:} IA do Google para analisar documentos
		\vspace{0.5em}
		\begin{block}{Como usar na educação inclusiva}
			\small
			• Fazer upload de PDFs e gerar resumos\\
			• Criar guias de estudo a partir de materiais\\
			• Gerar podcasts (Audio Overview) sobre conteúdos\\
			• Responder perguntas com base nos documentos
		\end{block}
		\textbf{Diferencial:} trabalha só com o que você envia
		\citabnt{Google (2024); AUGUSTO et al. (2026)}
	\end{frame}
	
	% ----- 5. Japiim (Bemol) -----
	\begin{frame}{\faIcon{map-marked-alt}~Japiim (Bemol)}
		\textbf{O que é:} plataforma de IA generativa da Bemol S/A, sediada em Manaus
		\vspace{0.5em}
		\begin{block}{Como usar na educação inclusiva}
			\small
			• Acesso a modelos da OpenAI, Groq, Anthropic, Google e Meta\\
			• Criação de Assistentes personalizados por tema\\
			• Atendimento em português com contexto amazônico\\
			• Apoio a estudantes ribeirinhos e indígenas
		\end{block}
		\textbf{Diferencial:} IA do Norte, democratiza acesso à IA no Brasil
		\citabnt{Japiim (2024); AUGUSTO et al. (2026)}
	\end{frame}
	
	% ----- 6. VLibras -----
	\begin{frame}{\faIcon{hands}~VLibras}
		\textbf{O que é:} tradutor automático para Libras (gov.br)
		\vspace{0.5em}
		\begin{block}{Como usar na educação inclusiva}
			\small
			• Instalar plugin em sites educacionais\\
			• Traduzir textos e vídeos para Libras\\
			• Tornar plataformas acessíveis para surdos\\
			• Usar em apresentações e aulas online
		\end{block}
		\textbf{Onde:} \url{https://vlibras.gov.br}
		\citabnt{Brasil (2024); AUGUSTO et al. (2026)}
	\end{frame}
	
	% ----- 7. Seeing AI -----
	\begin{frame}{\faIcon{eye}~Seeing AI}
		\textbf{O que é:} app da Microsoft que descreve o mundo
		\vspace{0.5em}
		\begin{block}{Como usar na educação inclusiva}
			\small
			• Descrever cenas e objetos para cegos\\
			• Ler textos em voz alta (OCR)\\
			• Identificar cores e pessoas\\
			• Apoiar mobilidade em espaços escolares
		\end{block}
		\textbf{Onde:} iOS e Android (gratuito)
		\citabnt{Microsoft (2024); AUGUSTO et al. (2026)}
	\end{frame}
	
	% ----- 8. Legendagem -----
	\begin{frame}{\faIcon{closed-captioning}~Legendagem automática}
		\textbf{O que é:} legendas em tempo real
		\vspace{0.5em}
		\begin{block}{Como usar na educação inclusiva}
			\small
			• Google Meet $\rightarrow$ Ativar legendas\\
			• YouTube $\rightarrow$ Legendas automáticas $\rightarrow$ editar\\
			• OBS Studio $\rightarrow$ legendas ao vivo\\
			• Aulas gravadas com transcrição automática
		\end{block}
		\textbf{Para quem:} surdos, baixa audição, não nativos
		\citabnt{Google (2024); AUGUSTO et al. (2026)}
	\end{frame}
	
	\begin{frame}{O papel do professor}
		\begin{center}
			\Large
			\textbf{IA não substitui. Amplia.}
		\end{center}
		\vspace{1em}
		\begin{center}
			Mediação humana é insubstituível.
		\end{center}
		\citabnt{AUGUSTO et al. (2026); VALENTE (2022); VYGOTSKY (2007)}
	\end{frame}
	
	\begin{frame}{Condições necessárias}
		\begin{center}
			\Large \textbf{Para funcionar:}
		\end{center}
		\vspace{1em}
		\begin{columns}[T]
			\begin{column}{0.5\textwidth}
				\begin{itemize}
					\item Formação continuada
					\item Infraestrutura
					\item Inclusão digital
				\end{itemize}
			\end{column}
			\begin{column}{0.5\textwidth}
				\begin{itemize}
					\item Ética
					\item Políticas públicas
					\item Transparência
				\end{itemize}
			\end{column}
		\end{columns}
		\citabnt{AUGUSTO et al. (2026)}
	\end{frame}
	
	\begin{frame}{Atividade prática — 30 min}
		\begin{center}
			\Large \textbf{Mãos na massa}
		\end{center}
		\vspace{1.5em}
		\begin{itemize}
			\item Escolha 1 ferramenta
			\item Teste com um caso real
			\item Compartilhe: funcionou? Por quê?
		\end{itemize}
	\end{frame}
	
	\begin{frame}{Ética e cuidado}
		\begin{center}
			\Large \textbf{Antes de usar:}
		\end{center}
		\vspace{1em}
		\begin{center}
			\Large
			Privacidade $\cdot$ Vieses $\cdot$ Transparência $\cdot$ Consentimento
		\end{center}
		\citabnt{COMISSÃO EUROPEIA (2019); UNESCO (2021)}
	\end{frame}
	
	% ============================================================
	% BLOCO 3 — O QUÊ
	% ============================================================
	\section{O QUÊ}
	
	\begin{frame}{}
		\begin{center}
			\Huge O QUÊ
			\vspace{1em}
			\normalsize
			O que é IA? O que aprendemos? O que fica?
		\end{center}
	\end{frame}
	
	\begin{frame}{Afinal, o que é IA?}
		\begin{center}
			\Large \textbf{Inteligência Artificial é:}
		\end{center}
		\vspace{1em}
		\begin{center}
			Sistemas que simulam capacidades humanas
			\vspace{0.5em}
			\normalsize
			aprender $\cdot$ raciocinar $\cdot$ perceber $\cdot$ decidir
		\end{center}
		\citabnt{RUSSELL; NORVIG (2022)}
	\end{frame}
	
	\begin{frame}{IA na educação é...}
		\begin{center}
			\Large \textbf{Não é mágica. É ferramenta.}
		\end{center}
		\vspace{1em}
		\begin{columns}[T]
			\begin{column}{0.5\textwidth}
				\centering
				\textbf{Adapta}\\
				\textbf{Apoia}\\
				\textbf{Amplia}
			\end{column}
			\begin{column}{0.5\textwidth}
				\centering
				\textbf{Mas não substitui}\\
				\vspace{0.5em}
				\small o professor
			\end{column}
		\end{columns}
		\citabnt{AUGUSTO et al. (2026)}
	\end{frame}
	
	\begin{frame}{O que os resultados mostram}
		\begin{center}
			\Large \textbf{Evidências:}
		\end{center}
		\vspace{1em}
		\begin{center}
			\Large
			$+$ Acessibilidade\\
			$+$ Personalização\\
			$+$ Autonomia\\
			$+$ Participação
		\end{center}
		\citabnt{AUGUSTO et al. (2026)}
	\end{frame}
	
	\begin{frame}{O que ainda desafia}
		\begin{center}
			\Large \textbf{Desafios:}
		\end{center}
		\vspace{1em}
		\begin{center}
			\Large
			Privacidade $\cdot$ Vieses $\cdot$ Formação $\cdot$ Desigualdade
		\end{center}
		\citabnt{AUGUSTO et al. (2026); UNESCO (2021)}
	\end{frame}
	
	\begin{frame}{O que fica}
		\begin{center}
			\Large \textbf{Três compromissos:}
		\end{center}
		\vspace{1.5em}
		\begin{enumerate}
			\item Incluir é propósito
			\item IA é meio, não fim
			\item Professor é insubstituível
		\end{enumerate}
		\citabnt{AUGUSTO et al. (2026)}
	\end{frame}
	
	\begin{frame}{Atividade final — 20 min}
		\begin{center}
			\Large \textbf{Mapa mental coletivo}
		\end{center}
		\vspace{1.5em}
		\begin{itemize}
			\item O que levo daqui?
			\item O que vou aplicar?
			\item O que preciso aprender mais?
		\end{itemize}
	\end{frame}
	
	\begin{frame}{Mensagem final}
		\begin{center}
			\Large
			``A IA não substitui o professor.\\
			\vspace{0.5em}
			Potencializa sua atuação.''
		\end{center}
		\citabnt{AUGUSTO et al. (2026)}
	\end{frame}
	
	\begin{frame}{}
		\begin{center}
			\Huge Obrigado.
			\vspace{1em}
			\normalsize
			Perguntas? Discussões?
			\vspace{2em}

			\texttt{2024007388@ifam.edu.br}
		\end{center}
	\end{frame}
	
	% ============================================================
	% REFERÊNCIAS
	% ============================================================
	\begin{frame}[allowframebreaks]{Referências}
		\footnotesize
		\begin{thebibliography}{9}
			
			\bibitem{augusto2026}
			AUGUSTO, Emerson Aparecido et al. \textbf{Inteligência artificial e educação inclusiva: transformações promovidas pelas tecnologias inteligentes nos processos de ensino e aprendizagem}. Revista Colombiana de Ciências e Humanidades (REHCOL), [S. l.], v. 1, n. 1, p. 1--13, 2026. DOI: 10.5281/zenodo.21418317.
			
			\bibitem{brasil2024}
			BRASIL. \textbf{VLibras: tradutor automático de Libras}. Brasília: Governo Federal, 2024. Disponível em: \url{https://vlibras.gov.br}. Acesso em: 29 set. 2026.
			
			\bibitem{comissao2019}
			COMISSÃO EUROPEIA. \textbf{Diretrizes éticas para uma inteligência artificial confiável}. Bruxelas: Comissão Europeia, 2019.
			
			\bibitem{deepseek2024}
			DEEPSEEK. \textbf{DeepSeek: AI for reasoning and code}. [S. l.]: DeepSeek, 2024.
			
			\bibitem{google2024}
			GOOGLE. \textbf{NotebookLM e Google Meet: recursos de acessibilidade}. Mountain View: Google, 2024.
			
			\bibitem{japiim2024}
			JAPIIM. \textbf{Plataforma de IA Generativa}. Manaus: Bemol S/A, 2024. Disponível em: \url{https://japiim.com.br}. Acesso em: 29 set. 2026.
			
			\bibitem{microsoft2024}
			MICROSOFT. \textbf{Seeing AI: um app gratuito para pessoas com deficiência visual}. Redmond: Microsoft, 2024.
			
			\bibitem{openai2024}
			OPENAI. \textbf{ChatGPT}. [S. l.]: OpenAI, 2024.
			
			\bibitem{perplexity2024}
			PERPLEXITY. \textbf{Perplexity AI}. [S. l.]: Perplexity, 2024.
			
			\bibitem{russell2022}
			RUSSELL, Stuart; NORVIG, Peter. \textbf{Inteligência Artificial}. 4. ed. Rio de Janeiro: LTC, 2022.
			
			\bibitem{unesco1994}
			ORGANIZAÇÃO DAS NAÇÕES UNIDAS PARA A EDUCAÇÃO, A CIÊNCIA E A CULTURA (UNESCO). \textbf{Declaração de Salamanca e linha de ação sobre necessidades educativas especiais}. Salamanca: UNESCO, 1994.
			
			\bibitem{unesco2021}
			ORGANIZAÇÃO DAS NAÇÕES UNIDAS PARA A EDUCAÇÃO, A CIÊNCIA E A CULTURA (UNESCO). \textbf{Recomendação sobre a ética da inteligência artificial}. Paris: UNESCO, 2021.
			
			\bibitem{valente2022}
			VALENTE, José Armando. Tecnologias digitais, inteligência artificial e educação: perspectivas para a aprendizagem. \textbf{Revista e-Curriculum}, São Paulo, v. 20, n. 2, p. 1--18, 2022.
			
			\bibitem{vygotsky2007}
			VYGOTSKY, Lev S. \textbf{A formação social da mente}. 7. ed. São Paulo: Martins Fontes, 2007.
			
		\end{thebibliography}
	\end{frame}
	
\end{document}